\section{Limitations}
\label{sec:limitations}

The empirical scope is one simulated greenhouse-tomato domain. The mechanistic crop model
improves traceability, but the economic coefficients, site distribution, factor ranges,
facility layout and observation channels are benchmark choices. Results do not establish
the same tool effects in chemistry, materials, biology, field agriculture or a physical
greenhouse.

The model study covers Luna and a smaller Sol replication through one provider interface,
one base prompt family and passive tool delivery. Provider outputs varied at temperature
zero, which is why the protocol treats generation seeds as within-site replicates. A null
effect for these aids does not show that statistical tools cannot help an agent; it bounds
the tested implementation and delivery mechanism.

Within-cycle observations are consequential only through the recommendation during the
current synchronous protocol. Units cannot be released, treatments cannot be modified, and
the next batch begins after terminal outcomes are available. The experiment therefore tests
whether inspection changes reasoning and endpoint recommendations, not the value of early
stopping, asynchronous replacement or adaptive treatment. Those actions require explicit
cleanup time, residual value and staggered-capacity semantics.

The common-history evaluator has a strong one-factor calibration but incomplete
higher-dimensional calibration. Its finite candidate search also supplies a feasible lower
bound on the world optimum. We consequently base primary comparisons on realized regret and
restrict posterior-risk decompositions to effects larger than measured evaluator error.

Finally, the Phase~3 scripted reference rows and Phase~4 LLM rows use independent held-out
site ranges. They provide common-scale descriptive anchors, but they are not paired
model-versus-baseline tests. Direct superiority claims require a jointly randomized matrix
on the same sites.
